\documentstyle[% 


righttag% 


,11pt% 


,amssymb% 


,russian%               remove in Eng. 


,izvru%                 Set izven% in Eng. 


]{amsart} 


 


% {=} {>} {\le}


\newcommand{\const}{\operatorname{const}} 


\newcommand{\sgn}{\operatorname{sgn}} 


\newcommand{\tts}[1]{} 


 


%\hoffset=-15mm%                  For Eng. 


\setcounter{page}{73} 


\god{2001} 


\nomer{10~(473)} %                        Remove in Eng. 


%\nomer{Vol.\,45, No.\,10}%               Set in Eng. 


%\str{\pageref{begin}--\pageref{end}}%  Set in Eng. 


\udk{517.95} 


 


\author{\it К.Б.\,Сабитов, А.А.\,Акимов} 


% NB:   ^^ remove in Eng. 


\title{К теории аналога задачи Неймана\\ для  


уравнений смешанного типа} 


\thanks{Работа выполнена при финансовой поддержке Российского фонда 


фундаментальных исследований, грант \No\,99-0100934, и Министерства


образования Российской Федерации 


по фундаментальным исследованиям в области математики, грант \No\,22 .} 


 


\date{} 


%\pagestyle{myheadings}%         Set in Eng. 


\pagestyle{plain} %              Remove in Eng. 


\begin{document} 


%\thispagestyle{plain}%          Set in Eng. 


\maketitle 


\label{begin} 


 


 


\section{Постановка задачи} 


 


Рассмотрим уравнение 


\begin{equation} 


L u=K(y)u_{xx}+u_{yy}+A u_x+B u_y+C u=F(x,y), 


\end{equation} 


где $y K(y)>0$ при $y\ne0$, в области $D$, ограниченной простой кривой 


$\Gamma$, лежащей в полуплоскости $y>0$ с концами в точках $A_1(0,0)$ и 


$B_1(l,0)$, $l>0$, и характеристиками $\gamma_1$ и $\gamma_2$ уравнения 


(1) при $y<0$ 


$$ 


\gamma_1:\xi=x+\int^y_0 \sqrt{-K(t)}\,dt=0,\quad 


\gamma_2:\eta=x-\int^y_0\sqrt{-K(t)}\,dt=l, 


$$ 


где $K(y)\in C[y_{c_1},0]\wedge C^2[y_{c_1},0]$, $y_{c_1}$ --- ордината 


точки $C_1$ пересечения 


характеристик $\gamma_1$ и $\gamma_2$. Пусть $D_+=D\cap\{y>0\}$, 


$D_-=D\cap\{y<0\}$. Для 


уравнения (1) в области $D$ поставим задачу типа Неймана, изученную 


Моравец [1]. 


\medskip 


 


{\bf Задача Моравец.} Найти функцию $u(x,y)$, удовлетворяющую условиям 


\begin{gather} 


u(x,y)\in C(\ov D)\wedge C^1(D\cup\Gamma)\wedge C^2(D_-\cup D_+);\\ 


L u(x,y)\equiv F(x,y),\quad (x,y)\in D_+\cup D_-; \\ 


K u_x dy-u_y dx=\phi(x,y),\quad (x,y)\in\Gamma; \\ 


K u_x dy-u_y dx=\psi(x,y),\quad (x,y)\in\gamma, 


\end{gather} 


где $\phi$ и $\psi$ --- заданные достаточно гладкие функции. 


\medskip 


 


Отметим, что задача (2)--(5) впервые изучена Моравец [1] для 


уравнения Чаплыгина (в этом случае в (1) $K(0)=0$, $K'(y)>0$ при 


всех $(x,y)\in\ov D$, $A(x,y)=B(x,y)=C(x,y)=F(x,y)=0$) 


в более общей 


постановке (в гиперболической части граничное условие (5) задано, 


вообще говоря, на неxарактеристической кривой) и в несколько иной 


области, чем $D$. Методом вспомогательных функций там доказана теорема 


единственности решения задачи при существенных геометрических ограничениях 


на кривую $\Gamma$ в точках подхода к оси $y=0$. 


В монографии Л.\,Берса ([2], с.\,118) для уравнения 


Лаврентьева--Бицадзе методом вспомогательных функций доказана 


теорема единственности решения задачи (2)--(5) в постановке 


Моравец, когда кривая $\Gamma$ в точках подхода к оси $y=0$ 


перпендикулярна, 


производная $\frac{dy}{dx}$ не обращается в нуль на $\Gamma$, за 


исключением одной 


точки. В случае, когда кривая $\Gamma$ совпадает с характеристикой, 


получена 


теорема единственности решения задачи Моравец для уравнения 


Чаплыгина [3] при условии $3K'{}^2-2K K''\ge0$, $y<0$, и 


Лаврентьева--Бицадзе 


[4] без каких-либо ограничений геометрического характера на 


кривую $\Gamma$. В данной работе при некоторых ограничениях на 


коэффициенты 


уравнения (1) установлены экстремальные свойства решений 


задачи (2)--(5) в области гиперболичности $D_-$ и смешанной области $D$. 


На основании этих свойств получены теоремы единственности решения 


задачи (2)--(5) без ограничений геометрического характера на кривую 


$\Gamma$. В п.\,4 приведены достаточные условия относительно функции 


$K(y)$, при которых справедливы теоремы единственности решения задачи 


Моравец для уравнения Чаплыгина без каких-либо ограничений 


геометрического характера на кривую $\Gamma$. 


 


\section{Экстремальные свойства решений задачи Моравец в 


области гиперболичности} 


 


В области $D_-$ перейдем к характеристическим координатам $(\xi,\eta)$. 


Тогда уравнение (1) примет вид 


\begin{equation} 


L_0u\equiv u_{\xi\eta}+a(\xi,\eta)u_\xi+b(\xi,\eta)u_\eta+ 


c(\xi,\eta)u=f(\xi,\eta), 


\end{equation} 


где 


\begin{gather*} 


a=\frac{(A+B\sqrt{-K}-K'/2\sqrt{-K})}{4K},\quad c=\frac{C}{4K},\\ 


b=\frac{(A-B\sqrt{-K}+K'/2\sqrt{-K})}{4K},\quad f=\frac{F}{4K}, 


\end{gather*} 


а область $D_-$ отображается в область $\Delta$, ограниченную отрезками 


$A_0B_0$ ($\eta=\xi$), $B_0C_0$ ($\eta=l$) и $A_0C_0$ ($\xi=0$). 


 


Пусть $\alpha=a\beta$, $\beta=\exp\displaystyle\int b\,d\xi$, функции 


$a(\xi,\eta)$, $a_\xi(\xi,\eta)$, $b(\xi,\eta)$, $c(\xi,\eta)$ 


непрерывны в $\ov\Delta$, кроме, быть может, отрезка $\ov{A_0B_0}$, и 


при $(\xi,\eta)\in\Delta\cup B_0C_0$ 


удовлетворяют одному из следующих условий: 


\begin{gather} 


\alpha(0,\eta)\ge0,\ \; h=\alpha_\xi+\alpha b-c\ge0,\ \; 


\int^\xi_0 \beta(t,\eta)c(t,\eta)dt>0, \ \; 0<\xi<\eta\le l; 


\tag{$B_1$} \\ 


% 


c(\xi,\eta)=0,\ \; \alpha(\xi,\eta)-\int^\xi_0 \beta(t,\eta) 


|h(t,\eta)|dt>0,\ \; 0<\xi<\eta\le l.\tag{$B_2$} 


\end{gather} 


Предполагается, что правая часть $f(\xi,\eta)$ уравнения (6) интегрируема 


по $\xi$ на каждом отрезке $[0,\xi_0]$ характеристики $\eta=\eta_0$, 


где $0<\xi_0<\eta_0\le l$. 


 


\begin{defn} 


Регулярным в $\Delta$ решением уравнения (6) назовем 


функцию $u(\xi,\eta)$, удовлетворяющую условиям: 


 


1) $u(\xi,\eta)\in C(\ov\Delta)\wedge C^1(\Delta)$, $u_{\xi\eta}\in 


C(\Delta)$, $L_0u\equiv0$ в $\Delta$; 


 


2) производная $u_\eta$ непрерывна в $\ov\Delta$, 


кроме, быть может, отрезка $\ov{A_0B_0}$. 


\end{defn} 


 


\begin{lem} 


Пусть $1)$ коэффициенты уравнения $(6)$ обладают отмеченной 


выше гладкостью и удовлетворяют условию $(B_1);$ $2)$~$f(\xi,\eta)\le0$ 


$(\ge0)$ на $\Delta\cup B_0C_0;$ $3)$~$u(\xi,\eta)$ --- регулярное в 


$\Delta$ решение уравнения 


$(6)$, удовлетворяющее условию $u_\eta=0$ на характеристике $A_0C_0$. 


Тогда, 


если $\max\limits_{\ov\Delta}(\xi,\eta)>0$


$(\min\limits_{\ov\Delta}(\xi,\eta)<0)$, то этот 


максимум $($минимум$)$ 


достигается на отрезке $\ov{A_0B_0}$. 


\end{lem} 


 


{\bf Доказательство.} Проинтегрируем тождество 


\begin{equation} 


\beta L_0(u)\equiv\frac{\partial}{\partial\xi} 


(\beta u_\eta+\alpha u)-\beta h u=f\beta 


\end{equation} 


пo отрезку $N M$ прямой $\eta=\const$, принадлежащему 


области $\Delta$. Тогда получим 


\begin{equation} 


(\beta u_\eta+\alpha u)\big|^M_N=\int_{N M}\beta h u\,d\xi+ 


\int_{N M}\beta f\,d\xi. 


\end{equation} 


В равенстве (8) отрезок $N M$ может принадлежать не только области 


$\Delta$, 


но и $\ov\Delta\setminus\ov{A_0B_0}$. Пусть 


$\max\limits_{\ov\Delta}u(\xi,\eta)=u(Q)>0$. Допустим, что 


$Q\notin\ov{A_0B_0}$. 


Тогда точка $Q\in\Delta\cup B_0C_0$. Из точки $Q$ проведем отрезок 


$\eta=\const$ 


до пересечения с характеристикой $\xi=0$ в точке $P$. В (8) в 


качестве отрезка $N M$ возьмем $P Q$. Тогда получим 


\begin{multline*} 


\beta(Q)u_\eta(Q)=\int_{P Q}\beta h u\,d\xi+ 


\int_{P Q}\beta f\,d\xi-\alpha(Q)u(Q)+\alpha(P)u(P)= \\ 


% 


=\int_{P Q}\beta f\,d\xi+\int_{P Q}\beta h[u-u(Q)]d\xi+ 


u(Q)\int_{P Q}\beta h\,d\xi-\alpha(Q)u(Q)+\alpha(P)u(P)= \\ 


% 


=\int_{P Q}\beta f\,d\xi+\int_{P Q}\beta h[u-u(Q)]d\xi-u(Q) 


\int_{P Q}\beta c\,d\xi-\alpha(P)[u(Q)-u(P)]. 


\end{multline*} 


Отсюда в силу условия $(B_1)$ из неравенств $f\le0$ в $\Delta\cup 


B_0C_0$ и $u(Q)>0$ следует 


$u_\eta(Q)<0$. Но это противоречит тому, что в точке $Q$ максимума 


функции $u(\xi,\eta)$ производная $u_\eta(Q)\ge0$. 


 


\begin{note} 


Лемма 1 по существу представляет собой результат 


Агмона, Ниренберга и Проттера [5], который приведен для удобства 


дальнейшего изложения. 


\end{note} 


 


\begin{lem} 


Пусть $1)$ коэффициенты уравнения $(6)$ обладают отмеченной 


выше гладкостью и удовлетворяют условию $(B_2);$ 


$2)$~$f(\xi,\eta)\equiv0$ 


нa $\Delta\cup B_0C_0;$ $3)$~$u(\xi,\eta)$ --- регулярное в $\Delta$ 


решение уравнения $(6)$, 


удовлетворяющее условию $u=0$ на характеристике $A_0C_0$. Тогда, если 


$\max\limits_{\ov\Delta}|u(\xi,\eta)|>0$, то этот максимум достигается на 


отрезке $\ov{A_0B_0}$. 


\end{lem} 


 


{\bf Доказательство.} По условию $u(\xi,\eta)$ непрерывна в $\ov\Delta$ 


и $u_\eta(0,\eta)=0$ 


на характеристике $A_0C_0$. Тогда $u(0,\eta)=\const=u(A_0)$ при всех 


$0\le\eta\le l$. 


Поскольку коэффициент $c(\xi,\eta)\equiv0$ в $\Delta$, то, не теряя общности, 


считаем $u=0$ на $A_0C_0$, т.\,к. в противном случае, рассматривая 


новую функцию $v(\xi,\eta)=u(\xi,\eta)-u(A_0)$, имели бы, что $v(\xi,\eta)$


в области $\Delta$ является регулярным решением уравнения (6) 


$$ 


L_0(v)\equiv v_{\xi\eta}+a v_\xi+b v_\eta=0, 


$$ 


$v_\eta=u_\eta=0$ на $A_0C_0$ и $v=0$ на $\ov{A_0C_0}$. Пусть 


$\max\limits_{\ov\Delta}|u(\xi,\eta)|=|u(Q)|>0$. 


B силу линейности и однородности уравнения (6) можно предположить, 


что $u(Q)>0$. Тогда $\max\limits_{\ov\Delta}u(\xi,\eta){=}u(Q){>}0$ и 


$|u(\xi,\eta)|\le u(Q)$ в $\ov\Delta$. 


Допустим, что точка $Q\in \Delta\cup B_0C_0$. Тогда, рассуждая аналогично 


доказательству 


леммы 1, из (8) получим 


\begin{equation} 


\beta(Q)u_\eta(Q)+\alpha(Q)u(Q)=\int_{P Q}\beta h u\,d\xi. 


\end{equation} 


Оценивая правую часть (9), имеем 


$$ 


\beta(Q)u_\eta(Q)\le\int_{P Q}\beta|h|\,|u|d\xi- 


\alpha(Q)u(Q)\le-u(Q)\bigg[\alpha(Q)-\int_{P Q} 


\beta|h|d\xi\bigg]. 


$$ 


Отсюда в силу условия $(B_2)$ следует $u_\eta(Q)<0$. Это противоречит 


тому, что $u_\eta(Q)\ge0$. Следовательно, наше допущение 


неверно, и $Q\in\ov{A_0B_0}$. 


 


\begin{note} 


Если $h=a_\xi+a b\ge0$ на множестве $\Delta\cup B_0C_0$, то 


условие $(B_2)$ равносильно системе неравенств $a(0,\eta)>0$, $h\ge0$, 


$c\equiv0$ 


при $0<\xi<\eta\le l$; если же $h=a_\xi+a b<0$ на множестве $\Delta\cup 


B_0C_0$, 


то условие $(B_2)$ равносильно системе неравенств 


$2\alpha(\xi,\eta)-\alpha(0,\eta)>0$, $h<0$, $c\equiv0$ 


при $0<\xi<\eta\le l$. 


\end{note} 


 


\begin{note} 


Если в леммах 1 и 2 краевое условие $u_\eta=0$ на левой 


характеристике $A_0C_0$ заменить на $u_\xi=0$ на правой характеристике 


$C_0B_0$, то для выполнения заключения этих лемм нужно заменить 


условия $(B_k)$, $k=1,2$, на следующие условия: 


\begin{gather} 


b(\xi,l)<0,\ \; h^*=b_\eta+a b-c\ge0,\ \; 


\int^l_\eta \beta^*(\xi,t)c(\xi,t)dt>0,\ \; 0\le\xi<\eta<l; 


\tag{$B^*_1$} \\ 


% 


c(\xi,\eta)=0,\ \; \alpha^*(\xi,\eta)+\int^l_\eta \beta^*(\xi,t) 


|h^*(\xi,t)|dt<0,\ \; 0\le\xi<\eta<l,\tag{$B^*_2$} 


\end{gather} 


где $\beta^*=\exp\displaystyle\int a\,d\eta$, $\alpha^*=b\beta^*$. 


Полученные таким образом утверждения назовем 


соответственно леммами $1^*$ и $2^*$. 


\end{note}


 


\section{Экстремальные свойства решений задачи Моравец в 


смешанной области} 


 


\begin{lem}[{\series{m}\shape{n}\selectfont [5], [6]}] 


Пусть $1)$ в области $D_+$ коэффициенты уравнения $(1)$ 


ограничены и $C(x,y)\le0;$ $2)$~$u(x,y)\in C(\ov D)\wedge C^1(D_+\cup A 


B)\wedge C^2(D_+)$, $L u\equiv F(x,y)\ge0$ $(\le0)$ в $D_+;$ 


$3)$~$\max\limits_{\ov{D_+}}u(x,y)=u(x_0,0)>0$ $(\min\limits_{\ov 


D_+}u(x,y)=u(x_0,0)<0)$, $0<x_0<l$. 


Тогда 


\begin{equation} 


\lim_{y\to0+0}u_y(x_0,y)=u_y(x_0,0+0)<0\ \; (>0). 


\end{equation} 


\end{lem} 


 


Обозначим через $\Omega$ достаточно малую окрестность точки $B_1$, и пусть 


$\Omega_1=D_+\cap\Omega$. 


 


\begin{lem} 


Пусть $1)$ выполнены условия $1)$ и $2)$ леммы $3;$ 


$2)$~коэффициенты уравнения $(1)$ и производные $A_x$ и $B_y$ непрерывны 


в $\ov{\Omega_1}$ 


и там удовлетворяют условиям $2C-A_x-B_y\le0$, $B(x,0)\ge0;$ 


$B\,dx-A\,dy\ge0$ на $\Gamma\cap\ov{\Omega_1};$ $3)$~$K u^2_x+u^2_y$ и 


$F(x,y)$ суммируемы в $\Omega_1$. Тогда, 


если $B_1$ является точкой изолированного положительного максимума 


$($отрицательного минимума$)$ функции $u(x,y)$, то существует точка 


$Q_1\in(\Gamma\cup A_1B_1)\cap\ov{\Omega_1}$ такая, что 


\begin{equation} 


\delta_s[u(Q_1)]>0\quad (\delta_s[u(Q_1)]<0). 


\end{equation} 


\end{lem} 





\begin{pf} 


Возьмем число $r\in(0,u(B_1))$, настолько близкое 


к числу $u(B_1)$, чтобы кривая $\Gamma_0$, составленная из линии уровня 


$u(x,y)=r$, целиком лежала в малой окрестности $\Omega$ точки $B_1$ и при 


всех $(x,y)$, принадлежащих области $G$, ограниченной кривыми $B_1B_2$, 


$\Gamma_0$ и отрезком $B_0B_1$, $u(x,y)\ge r$. Здесь $B_0$ и $B_2$ --- 


точки пересечения 


кривой $\Gamma_0$ соответственно с отрезком $A_1B_1$ и кривой $\Gamma$. 


Существование 


такой кривой $\Gamma_0$ обосновывается аналогично работе [7], причем 


кривая $\Gamma_0$ является кусочно-гладкой. В области $G$ рассмотрим 


функцию 


$v(x,y)=u(x,y)-r$, которая является решением уравнения 


\begin{gather} 


L v(x,y)\equiv K(y)v_{xx}+v_{yy}+A v_x+B v_y+C v=F-C r\ge0,\\ 


% 


v(x,y)=0\ \; \text{на}\ \; \Gamma_0. 


\end{gather} 


Предположим, что неравенство (11) не выполняется, т.\,е. 


\begin{equation} 


\delta_s[u(x,y)]\le0\ \; \forall(x,y)\in B_0B_1\cup B_1B_2. 


\tag{$13^*$} 


\end{equation} 


Интегрируя тождество 


$$ 


v L v=(K v v_x+\tfrac{1}{2}A v^2)_x+(v v_y+\tfrac{1}{2}B v^2)_y- 


K v^2_x-v^2_y+v^2(C-\tfrac{1}{2}A_x-\tfrac{1}{2}B_y) 


$$ 


по области $G$ с учетом условий (12), (13) и ($13^*$), получим 


\begin{multline} 


\int_{B_0B_1}[2v(x,0)v_y(x,0)+B(x,0)v^2(x,0)]dx+ 


\int_{B_1B_2}v^2(B\,dx-A\,dy)+ \\ 


% 


+\int_{B_1B_2}v(v_y\,dx-K v_x\,dy)+\iint_G (K v^2_x+v^2_y)dx\,dy- \\ 


% 


-\iint_G[v^2(2C-A_x-B_y)+2(C r-F)v]dx\,dy=0. 


\end{multline} 


Отсюда следует $v=0$, т.\,е. $u=\const$ в $\ov{G_+}$, а это противоречит 


тому, что $B_1$ является точкой изолированного максимума.


\end{pf}


 


\begin{defn} 


Регулярным в области $D$ решением уравнения (1) 


назовем функцию, удовлетворяющую условиям (2) и (3), и дополнительно 


потребуем, чтобы производная $u_\eta$ была непрерывной в замкнутой 


области $\ov{D_-}$, кроме, быть может, отрезка $\ov{A_1B_1}$. 


\end{defn} 


 


\begin{thm} 


Пусть $1)$ коэффициенты уравнения $(1)$ в области $D_+$ 


ограничены и $C(x,y)\le0;$ $2)$~коэффициенты уравнения $(1)$ в области 


$D_-$ в характеристических координатах $(\xi,\eta)$ удовлетворяют условиям 


леммы~$1;$ $3)$~$F(x,y)\ge0$ $(\le0)$ на $D_+\cup D_-;$ $4)$~$u(x,y)$ --- 


регулярное в 


$D$ решение уравнения $(1)$, удовлетворяющее условию $(5)$, где функция 


$\psi(x,y)=0$ на характеристике $\gamma_1;$ $5)$~$\max\limits_{\ov 


D}u(x,y)>0$ $(\min\limits_{\ov D}u(x,y)<0)$. 


Тогда $\max\limits_{\ov D}u(x,y)$ $\big(\min\limits_{\ov D}u(x,y)\big)$ 


достигается на кривой $\ov\Gamma$. 


\end{thm} 


 


{\bf Доказательство.} Пусть $\max\limits_{\ov D}u(x,y)=u(Q)>0$. В силу 


внутреннего 


принципа экстремума для эллиптических уравнений $u(x,y)\ne\const$ 


не может принимать положительного максимума $u(Q)$ в области $D_+$. 


Поэтому значение $u(Q)$ принимается на $\ov{D_-}$ или на $\ov\Gamma$. 


Допустим, что 


$Q\in\ov{D_-}$. По условию~4) теоремы $K u_x\,dy-u_y\,dx=0$ на 


характеристике 


$\gamma_1$: $\sqrt{(-K)}dy+dx=0$. Отсюда следует, что вдоль кривой 


$\gamma_1$ производная 


$u_\eta=0$. Тогда в этом случае на основании леммы~1 точка 


$Q\in\ov{A_1B_1}$. 


Пусть $Q\in A_1B_1$, т.\,е. $Q=(x_0,0)$, $0<x_0<l$. В этой точке из 


гиперболической 


части области $D$ имеем $u_y(x_0,-0)\ge0$, что в силу леммы~3 


противоречит неравенству (10). Следовательно, $Q\in\ov\Gamma$. 


 


\begin{thm} 


Пусть $1)$ коэффициенты уравнения $(1)$ в области $D_+$ 


ограничены и $C(x,y){\le}0;$ $2)$~коэффициенты уравнения $(1)$ в области 


$D_-$ в характеристических координатах $(\xi,\eta)$ удовлетворяют условиям 


леммы~$2;$ $3)$~$F(x,y)\equiv0$ в~$D;$ $4)$~$u(x,y)$ --- регулярное в $D$ 


решение 


уравнения $(1)$, удовлетворяющее условию $(5)$, где функция 


$\psi(x,y)\equiv0$. 


Тогда, если $\max\limits_{\ov D}|u(x,y)|{>}0$, то этот максимум 


достигается на кривой $\ov\Gamma$. 


\end{thm} 


 


{\bf Доказательство.} Поскольку $K u_x\,dy-u_y\,dx=0$ на характеристике 


$\gamma_1$, 


то $u_\eta=0$ вдоль этой кривой. Пусть $\max\limits_{\ov 


D}|u(x,y)|=|u(Q)|>0$. В 


силу линейности и однородности уравнения (1) можно полагать, что 


$u(Q)>0$. Рассуждая аналогично доказательству теоремы 1, получим 


$Q\in\ov{D_-}\cup\ov\Gamma$. Допустим, что $Q\in\ov{D_-}$. Тогда по 


лемме 2 точка $Q\in\ov{A_1B_1}$. 


Пусть $Q\in A_1B_1$, т.\,е. $Q=(x_0,0)$, $0<x_0<l$. В этой точке из 


гиперболической части $D_-$ имеем $u_y(x_0,-0)\ge0$, что на основании 


леммы 


3 противоречит неравенству (10). Значит, $Q\in\ov\Gamma$. 


 


\begin{thm} 


Пусть $1)$ выполнены условия леммы $4$ и условие $2)$ 


теоремы $1;$ $2)$~$u(x,y)$ --- решение однородной задачи $(2)$--$(5)$ из 


класса регулярных в $D$ решений уравнения $(1);$ $3)$ $u(0,0){=}0$. Тогда 


$u(x,y)\equiv0$ в $\ov D$. 


\end{thm} 


 


{\bf Доказательство.} Допустим, что существует точка $(x_1,y_1)\in D$ 


такая, 


что $u(x_1,y_1)\ne0$, $u(x,y)\ne\const$ в $D_+$. Пусть для определенности 


$u(x_1,y_1)>0$. Тогда $\max\limits_{\ov D}u(x,y)=u(Q){>}0$, $Q\in\ov D$. В 


силу 


теоремы 1 точка $Q\in\Gamma\cup B_1$. Пусть $Q\in \Gamma$. В этом 


случае на основании 


граничного принципа экстремума для эллиптических уравнений 


в точке $Q$ имеем неравенство $\delta_s(u)=K 


u_x\frac{dy}{ds}-u_y\frac{dx}{ds}\big|_Q>0$, которое противоречит 


равенству $\delta_s(u)=0$ на $\Gamma$. Пусть $Q=B_1$. В этом случае 


$B_1$ 


является точкой изолированного положительного глобального максимума 


функции $u(x,y)$. Тогда линии уровня $u(x,y)=r$ функции $u(x,y)$ 


в $D_+$, где $r\in(0,u(Q))$, при $r$, близких к числу $u(Q)$, будут 


располагаться 


концентрично вокруг точки $B_1$ с концами на кривой 


$A_1B_1\cup\Gamma$ [8]. Отсюда следует, что функция $u(x,0)$ в малой 


окрестности 


$B_1$, т.\,е. на $[b_0,l)$ оси $y=0$ возрастает. В противном случае 


существуют такие $x_1,x_2\in[b_0,l)$, что $x_1<x_2$ и $u(x_1,0)=u(x_2,0)$. 


Пусть $u(x_1,0)=u(x_2,0)$. Рассмотрим линии уровня $u(x,y)=u(x_1,0)$, 


$u(x,y)=u(x_2,0)$. Функция $u(x,0)$ на отрезке $[x_1,x_2]$ имеет 


экстремум, 


который достигается в некоторой точке $(x_0,0)$, $x_0\in(x_1,x_2)$. 


Линия уровня $u(x,y)=u(x_0,0)$ будет расположена между линиями 


$u(x_1,y)=u(x_1,0)$ и $u(x,y)=u(x_2,0)$. Это означает, что на линии 


$u(x,y)=u(x_0,0)$ функция $u(x,y)$ в области $D_0$ имеет локальный 


экстремум, 


чего быть не может. Если $u(x_1,y)>u(x_2,0)$, то существует 


$x_3\in[x_2,l)$ такая, что $u(x_1,0)=u(x_3,0)$. Тогда, рассуждая 


аналогично, 


снова получим противоречие. Пусть $\xi$ --- любая точка из $[b_0,l)$. Из 


точки 


$E(\xi,0)$ опустим перпендикуляр с основанием в $Q_2\in D_-$ и из этой 


точки 


проведем характеристику уравнения (1) до пересечения с 


характеристикой $A_1C_1$ в точке $C_0$. Область, ограниченную кривыми 


$A_1C_0$, $C_0Q_2$, 


$Q_2E$ и $E A_1$, обозначим через $H$. На основании леммы 1 можно 


показать, 


что $\max\limits_{\ov H}u(x,y)>0$ достигается на отрезке $\ov{A_1E}$ в 


точке $E$. Значит, 


$u_y(\xi,0-0)\ge0$ при всех $\xi\in[b_0,l)$. В силу леммы 4 в любой 


достаточно 


малой окрестности $\Omega$ точки $B_1$ на $\Gamma\cup A_1B_1$ существует 


точка $Q_1$ такая, 


что $\delta_s[u(Q_1)]>0$. Если $Q_1\in\ov{\Omega_1}\cap\Gamma$, то получим 


противоречие с граничным 


условием $\delta_s[u(Q_1)]=0$ на $\Gamma$. Пусть 


$Q_1\in\ov{\Omega_1}\cap A_1B_1$. Тогда на 


отрезке $[b_0,l)$ существует точка $(\xi_0,0)$ такая, что 


$u_y(\xi_0,0+0)<0$, а это 


противоречит тому, что $u_y(\xi_0,0-0)\ge0$. Следовательно, 


$u(x,y)\equiv\const$. 


С учетом условия 3)\; $u(x,y)\equiv0$ в $\ov D$. 


 


\begin{thm} 


Пусть $1)$ выполнены условия $1)$ и $2)$ теоремы $2$ и условия 


леммы $4;$ $2)$~$C(x,y){\equiv}0$ в области $D_+$. Тогда $u(x,y)=\const$ в 


$D$. \tts{remove {} around \equiv}


\end{thm} 


 


{\bf Доказательство.} В области $D$ введем вспомогательную функцию 


$W(x,y)=u(x,y)-u(A_1)$, 


которая является решением однородной задачи (2)--(5), 


где $\phi(x,y)\equiv0$, $\psi(x,y)\equiv0$, $F(x,y)\equiv0$, 


$C(x,y)\equiv0$, причем $W(A_1)=0$. 


Докажем, что $W(x,y)\equiv0$ в $D$. Допустим, что $W(x,y)\ne\const$ в 


области $D_+$. 


Тогда $\max\limits_{\ov D}|W(x,y)|=|W(Q)|>0$. В силу теоремы 2 точка 


$Q\in\Gamma\cup B_1$. 


Не теряя общности рассуждений, считаем, что 


$W(Q)>0$. Далее, рассуждая аналогично доказательству теоремы 3, 


получим $W(x,y)\equiv0$ в $D$. Отсюда следует $u(x,y)\equiv u(A_1)$ в 


$D$. 


 


В качестве примера рассмотрим уравнение 


\begin{equation} 


\sgn y|y|^n u_{xx}+u_{yy}+a_0|y|^{\frac{n-1}{2}}u_x=0, 


\end{equation} 


где $n>0$, $a_0=\const$, которое изучалось многими авторами [8], [9]. В 


характеристических координатах $(\xi,\eta)$ уравнение (15) принимает вид 


\begin{equation} 


u_{\xi\eta}+\frac{p_1}{n-\xi}u_\xi-\frac{p_2}{n-\xi}u_\eta=0, 


\end{equation} 


где 


$$ 


p_1=\frac{n-2a_0}{2(n+2)},\quad p_2=\frac{n+2a_0}{2(n+2)}. 


$$ 


В случае уравнения (16) 


\begin{gather*} 


a=\frac{p_1}{n-\xi},\quad b=-\frac{p_2}{n-\xi},\quad 


c=0,\ \; \beta=(n-\xi)^{p_2},\\ 


\alpha=p_1(n-\xi)^{p_2-1},\quad h=\frac{p_1(1-p_2)}{(n-\xi)^2}. 


\end{gather*} 


Отсюда нетрудно видеть, что коэффициенты $a$, $b$ и $a_\xi$ непрерывны в 


$\ov\Delta\setminus\ov{A_0B_0}$ и при $a_0<n/2$ удовлетворяют условию 


$(B_2)$. Следовательно, 


если $a_0<n/2$, то для уравнения (14) справедлива теорема 2 и все 


утверждения, 


вытекающие из этой теоремы. Теорема 4 о единственности решения 


задачи (2)--(5) для уравнения (14) имеет место при $a_0\le0$. 


Последнее ограничение вызвано условиями (11). 


 


\section{Единственность решения задачи Моравец для уравнения 


Чаплыгина} 


 


Рассмотрим в области $D$ уравнение Чаплыгина 


\begin{equation} 


K(y)u_{xx}+u_{yy}=0, 


\end{equation} 


где $K(0)=0$, $K'(y)>0$ при $y\in[y_{c_1},y_m]$, 


$y_m=\max\{y\mid(x,y)\in\ov D\}$, 


$K(y)$ --- достаточно гладкая на $[y_{c_1},y_m]$ функция. Именно 


для этого уравнения впервые была поставлена задача (2)--(5) в газовой 


динамике [11], [2]. 


 


В характеристических координатах уравнение (17) имеет вид (6), где 


\begin{gather*} 


a=\frac{1}{8}K'(y)(-K)^{-\frac{3}{2}},\quad 


b=-\frac{1}{8}K'(y)(-K)^{-\frac{3}{2}},\quad c=f\equiv0,\\ 


% 


\beta=(-K)^{\frac{1}{4}},\quad \alpha=a\beta=\frac{1}{8} 


K'(y)(-K)^{-\frac{5}{4}}, \\ 


% 


h=a_\xi+a b=a'_y y'_\xi+a b=\frac{1}{64(-K)^3}[5K'{}^2-4K K'']. 


\end{gather*} 


Если $5K'{}^2-4K K''\ge0$ в $D_-\cup C_1B_1$, то коэффициенты уравнения 


(17) 


удовлетворяют условию $(B_2)$. Если же $5K'{}^2-4K K''<0$ в $D_-\cup 


C_1B_1$, 


то при условии 


$$ 


2\alpha(\xi,\eta)-\alpha(0,\eta)=\frac{1}{8}K'(y)(-K)^{-\frac{5}{4}} 


\bigg[2-\frac{K'_0(y)}{K'(y)}\bigg(\frac{K}{K_0}\bigg)^{\frac{5}{4}} 


\bigg]>0, 


$$ 


где $K_0(y)=K(y)\big|_{\xi=0}$, $K'_0(y)=K'(y)\big|_{\xi=0}$, также 


выполнено условие 


$(B_2)$ (см. замечание 2). 


 


Таким образом, справедлива 


 


\begin{thm} 


Если $5K'{}^2\ge4K K''$ или $5K'{}^2<4K K''$ и 


$2-\frac{K'_0(y)}{K'(y)}\big(\frac{K}{K_0}\big)^{\frac{5}{4}}>0$ 


при $y<0$, то для уравнения $(17)$ справедлива теорема $2$, в частности, 


теорема $4$ о единственности решения задачи $(2)$--$(5)$. 


\end{thm} 


 


Теперь на основании результатов Проттера [12], [13] можно получить 


другие достаточные условия единственности решения задачи (2)--(5) для 


уравнения Чаплыгина. 


 


\begin{thm} 


Пусть $1)$ выполнено одно из следующих условий\/$:$ 


{\em а)}~функция $F(y)=2\big(\frac{K}{K'}\big)'+1=\frac{1}{K'{}^2} 


(3K'{}^2-2K K'')\ge0$ при $y<0;$ {\em б)}~существует $d_1=\const<0$, 


зависящая от размеров области $D_-$, такая, 


что $F(y)\ge d_1;$ {\em в)}~существует $d_2=\const>0$ такая, что 


$y_m<d_2$, a 


функция $F(y)$ при $y<0$ может принимать любые конечные значения\/$;$ 


$2)$~$u(x,y)$ --- решение однородной задачи $(2)$--$(5)$ для уравнения 


$(17)$, 


при этом $K u^2_x+u^2_y$ суммируема в малых окрестностях точек $A_1$ и 


$B_1$. 


 


Тогда $u(x,y)\equiv\const$ в $D$. 


\end{thm} 


 


{\bf Доказательство.} Введем вспомогательную функцию 


$v(x,y)=u(x,y)-u(C_1)$, 


где $u(C_1)$ --- значение функции $u(x,y)$ в точке $C_1$, которая является 


решением однородной задачи (2)--(5) для уравнения (17), 


причем $v(C_1)=0$. Из условия $K v_x\,dy-v_y\,dx=0$ на характеристике 


$A_1C_1$ 


что равносильно $v_\eta(0,\eta)=0$ при $0<\eta<l$, и $v(C_1)=0$ вытекает 


$v(x,y)=0$ на $\ov{A_1C_1}$. Далее, проведя рассуждения, аналогичные 


доказательству единственности решения задачи Трикоми для 


уравнения Чаплыгина [12], [13], [8], получим $v(x,y)=0$ в $D$. Отсюда 


следует $u(x,y)\equiv u(C_1)$ в $D$. 


 


\begin{thebibliography}{99} 


 


\bibitem{1} 


Morawetz C.S. {\em Note on a maximum principle and a uniqueness theorem 


for an elliptic-hyperbolic equation} // Proc. Roy. Soc. -- 1956. -- 


V.\,236. -- \No\,1024. -- С.\,141--144. 


\bibitem{2} 


Берс Л. {\em Математические вопросы дозвуковой и околозвуковой газовой 


динамики}. -- М.: Ин. лит., 1961. -- 208 с. 


\bibitem{3} 


Сабитов К.Б. {\em О задаче Моравец для уравнений смешанного типа} // 


Международн. семин. ``Дифференциальные уравнения и их 


приложения''. -- Самара, 27--30 июня 1995. -- С.\,73. 


\bibitem{4} 


Сомова Н.В. {\em О задаче Неймана--Моравец для уравнения 


Лаврентьева--Бицадзе} // Международн. семин. ``Дифференциальные уравнения 


и их приложения''. -- Самара, 27--30 июня 1995. -- С.\,78. 


\bibitem{5} 


Agmon S., Nirenberg L., Protter M.N. {\em A maximum principle for a 


class of hyperbolic equations and applications to equations of mixed 


elliptic-hyperbolic type} // 


Comm. Appl. Math. -- 1953. -- V.\,6. -- \No\,4. -- P.\,455--470. 


\bibitem{6} 


Сабитов К.Б. {\em К вопросу о существовании решения задачи Трикоми} // 


Дифференц. уравнения. -- 1992. -- Т.\,28. -- \No\,12. -- С.\,2092--2101. 


\bibitem{7} 


Сабитов К.Б., Капустин Н.Ю. {\em О решении одной проблемы в теории 


задачи Франкля для уравнений смешанного типа} // Дифференц. 


уравнения. -- 1991. -- Т.\,27. -- С.\,60--68. 


\bibitem{8} 


Смирнов М.М. {\em Уравнения смешанного типа}. Учеб. пособие. -- М.: 


Высш. школа, 1985. -- 304~с. 


\bibitem{9} 


Бицадзе А.В. {\em Некоторые классы уравнений в частных производных}. -- 


М.: Наука, 1981. -- 448~с. 


\bibitem{10} 


Смирнов М.М. {\em Вырождающиеся эллиптические и гиперболические 


уравнения}. -- М.: Наука, 1966. -- 292 с. 


\bibitem{11} 


Франкль Ф.И. {\em Избранные труды по газовой динамике}. -- М.: 


Наука, 1973. -- 712 с. 


\bibitem{12} 


Protter М.H. {\em Uniqueness theorems for Trikomi problem}. I // J. 


Rat. Mech. and Anal. -- 1953. -- V.\,2. -- \No\,1. -- Р.\,107--114. 


\bibitem{13} 


Protter М.H. {\em Uniqueness theorems for Trikomi problem}. II // J. 


Rat. Mech. and Anal. -- 1955. -- V.\,4. -- \No\,5. -- Р.\,721--733. 


 


\end{thebibliography} 


 


\vskip 2cm 


 


\post{\it Стерлитамакский государственный}{\it Поступила} 


\post{\it педагогический институт}{12.11.1999} 


\smallskip 


\post{\it Стерлитамакский филиал}{} 


\post{\it АН Республики Башкортостан}{} 


 


\label{end} 


\end{document} 


